%%=============================================================================
%% Methodologie
%%=============================================================================

\chapter{\IfLanguageName{dutch}{Methodologie}{Methodology}}%
\label{ch:methodologie}

%% TODO: In dit hoofstuk geef je een korte toelichting over hoe je te werk bent
%% gegaan. Verdeel je onderzoek in grote fasen, en licht in elke fase toe wat
%% de doelstelling was, welke deliverables daar uit gekomen zijn, en welke
%% onderzoeksmethoden je daarbij toegepast hebt. Verantwoord waarom je
%% op deze manier te werk gegaan bent.
%% 
%% Voorbeelden van zulke fasen zijn: literatuurstudie, opstellen van een
%% requirements-analyse, opstellen long-list (bij vergelijkende studie),
%% selectie van geschikte tools (bij vergelijkende studie, "short-list"),
%% opzetten testopstelling/PoC, uitvoeren testen en verzamelen
%% van resultaten, analyse van resultaten, ...
%%
%% !!!!! LET OP !!!!!
%%
%% Het is uitdrukkelijk NIET de bedoeling dat je het grootste deel van de corpus
%% van je graduaatsproef in dit hoofstuk verwerkt! Dit hoofdstuk is eerder een
%% kort overzicht van je plan van aanpak.
%%
%% Maak voor elke fase (behalve het literatuuronderzoek) een NIEUW HOOFDSTUK aan
%% en geef het een gepaste titel.

Het onderzoek verliep in vier opeenvolgende fasen.

\section{Fase 1: Literatuurstudie}%
\label{sec:fase-literatuurstudie}

In de eerste fase werd de nodige achtergrondkennis opgebouwd over API gateways, observability-pijlers en het OpenTelemetry-ecosysteem. Concreet werd de APISIX-documentatie grondig doorgenomen, met bijzondere aandacht voor de beschikbare observability-plugins (\texttt{prometheus}, \texttt{opentelemetry}, \texttt{zipkin}, \texttt{http-logger}) en hun configuratiemogelijkheden~\autocite{APISIX2024}. Daarnaast werd de officiële ``Observe Your API''-tutorial geëvalueerd~\autocite{APISIXObserveTutorial2024} en werden de documentaties van de kandidaat-opslagcomponenten (Thanos, Tempo, Loki, Alloy, OTEL Collector) bestudeerd.

Deliverable: een overzicht van de beschikbare plugins en integratiepatronen, en een antwoord op de vraag of APISIX OTLP-metrics ondersteunt.

\section{Fase 2: Architectuurontwerp}%
\label{sec:fase-architectuurontwerp}

Op basis van de literatuurstudie werd een architectuur uitgetekend voor de proof of concept. De centrale ontwerpkeuze was de rol van de OTEL Collector als centrale hub: zowel het Prometheus-scrape-eindpunt van APISIX als de OTLP-tracesstroom worden via de Collector gerouteerd naar hun respectieve backends. Dit patroon koppelt de gateway los van de specifieke opslagbackends.

De keuze voor de individuele componenten werd als volgt verantwoord:

\begin{itemize}
  \item \textbf{Thanos} voor metrics: langetermijn-retentie via \texttt{remote\_write}, compatibel met de \texttt{prometheusremotewrite} exporter van de OTEL Collector.
  \item \textbf{Tempo} voor traces: OTLP-native, lage opslagkost, integreert met Grafana.
  \item \textbf{Loki} voor logs: efficiënte label-gebaseerde indexering, LogQL-querytaal, native trace-ID-koppeling in Grafana.
  \item \textbf{Grafana Alloy} voor log-collectie: vervangt het deprecated Promtail, Docker-socket-discovery out-of-the-box.
  \item \textbf{Mock-upstream-services in Go}: drie eenvoudige HTTP-services die gestructureerde JSON-logs schrijven naar stdout, zodat de log-aggregatiepijpeline getest kan worden zonder afhankelijkheid van echte Be-Mobile-backends.
\end{itemize}

Deliverable: een architectuurdiagram en een overzicht van de te configureren componenten.

\section{Fase 3: Implementatie van de POC}%
\label{sec:fase-implementatie}

De derde fase bestond uit het opzetten van de volledige stack als Docker Compose-omgeving. De implementatie verliep iteratief: eerst werd de metrics-pijler werkend gemaakt (APISIX Prometheus-plugin → OTEL Collector → Thanos), daarna de traces-pijler (OpenTelemetry-plugin → OTEL Collector → Tempo), en tot slot de logs-pijler (mock-services stdout → Alloy → Loki).

Per component werden de configuratiebestanden als code beheerd in een Git-repository. Tijdens de implementatie werden configuratievalkuilen geïdentificeerd en gedocumenteerd.

Deliverable: een werkende, reproduceerbare Docker Compose-omgeving met alle configuratiebestanden in versiebeheer.

\section{Fase 4: Validatie en dashboards}%
\label{sec:fase-validatie}

In de laatste fase werden de Grafana-dashboards ontworpen en geïmplementeerd op basis van de beschikbare metrics, log-queries en trace-links. De stack werd gevalideerd door testverzoeken te sturen via een \texttt{make seed}-commando dat 60 requests verdeelt over alle geconfigureerde consumers, en door de resulterende data in alle drie pijlers te controleren.

Deliverable: zes functionele Grafana-dashboards, een validatierapport van de POC, en een architectuuraanbeveling voor productie.

